\documentclass[11pt]{article}
\usepackage[margin=1in]{geometry}
\usepackage{amsmath,amssymb,amsthm,mathtools}
\usepackage{array,booktabs,longtable}
\usepackage[hyphens]{url}
\usepackage[hidelinks]{hyperref}
\newcommand{\Q}{\mathbb Q}
\newcommand{\F}{\mathbb F}
\newcommand{\SO}{\operatorname{SO}}
\newcommand{\tr}{\operatorname{tr}}
\newcommand{\adj}{\operatorname{adj}}
\newcommand{\Span}{\operatorname{span}}
\newtheorem{theorem}{Theorem}[section]
\newtheorem{proposition}[theorem]{Proposition}
\newtheorem{corollary}[theorem]{Corollary}
\newtheorem{lemma}[theorem]{Lemma}
\theoremstyle{remark}
\newtheorem{remark}[theorem]{Remark}
\setlength{\emergencystretch}{3em}
\title{Issues and Resolutions to Open Problems in\\
  \emph{Rational Rotations of Three Dimensional Euclidean Space}}
\author{}
\date{12 October 2026}

\begin{document}
\maketitle
\begin{abstract}
This companion records the mathematical findings of the constructive Lean
formalisation of N.~J.~Wildberger's revised Euclidean rotation paper.
It gives an explicit repair to the
perpendicular-direction argument for null anchors, answers the
associativity question in Problem~80, gives exact finite-field counts and
conjugacy types for Problem~81, and states the finite-word and finite-order
results obtained for Problem~82. It also supplies the determinant-weighted
Clifford correspondence requested in Problem~83. Rational normalization,
positive rays, projective proportions and finite formulations of size claims
are kept distinct. Numbering refers to the revised source, not its earlier
August draft. The remaining subgroup-classification and arithmetic questions
are identified explicitly.
\end{abstract}

\section{Source, conventions and status}

The source is the supplied 33-page revision of \emph{Rational rotations of
three dimensional Euclidean space} \cite{source}. Its numbered assertion
inventory ends at Example~79, followed by Problems~80--83. We use row vectors,
\[
 B(v,u)=v\cdot u,\qquad Q(v)=v\cdot v,
 \qquad p(\rho)=\frac{3-\tr\rho}{4}.
\]
All general-field statements below assume characteristic different from two.
An assertion involving an ordered field additionally assumes its ordinary
ordered-field laws. The original quaternion product has the row-action signs
\begin{equation}
 (x:v)(y:u)=(xy-v\cdot u:xu+yv-v\times u),
 \qquad N(x:v)=x^2+Q(v).
 \label{eq:quaternion}
\end{equation}
The displayed representatives, their nonzero scales and their exact
coefficients are retained throughout.

The revised paper already repairs several defects discussed for the August
draft. In particular, Corollary~51 now requires both input caps to be defined,
Theorem~54 states the cap domains needed for its proof of conditional
associativity, and Problem~80 has fixed linear coefficients and excludes the
zero cross coefficient. The old projection counterexample to the earlier
four-coefficient conjecture is consequently not a counterexample to the new
question. It should not be carried forward as one.

\begin{center}
\begin{tabular}{>{\raggedright\arraybackslash}p{0.16\textwidth}>{\raggedright\arraybackslash}p{0.75\textwidth}}
\toprule
Source item & Result of this companion \\
\midrule
Theorem~33 & The formula is retained; a missing zero/null-anchor existence
case in its printed proof is supplied. \\
Problem~80 & Complete classification in characteristic different from two:
\(\delta=-\gamma^2\), using explicit scalar probes. \\
Problem~81 & Exact counts, spinor classes and conjugacy types; one
nonidentity parabolic conjugacy class in the full rotation group. \\
Problem~82 & Exact finite-word spinor image, kernel/coset tests and
finite-order table. Internal classification of generated subgroups remains
open. \\
Problem~83 & The projective Clifford correspondence for every displayed
nonsingular symmetric ternary form. \\
\bottomrule
\end{tabular}
\end{center}

This is an account of proved statements and their limits, rather than a
claim to have solved every research question in the source. Scope boundaries
concerning completed infinite objects and topology are not mathematical
errors in the source.

\section{Theorem 33: non-null perpendicular directions}

Definition~32 permits a cap with a zero or null anchor. The last sentence of
the proof of Theorem~33 appeals to nondegeneracy of the plane perpendicular
to a \emph{non-null} anchor. This establishes existence in that case, but
does not cover the identity anchor or the nonzero null anchor in Example~45.
The spread formula itself remains correct.

\begin{lemma}[Four-candidate perpendicular construction]
For every \(v=(t,r,s)\in F^3\), one can explicitly find \(w\) with
\(w\cdot v=0\) and \(Q(w)\ne0\).
\end{lemma}
\begin{proof}
Test the three perpendicular vectors
\[
 w_1=(-r,t,0),\qquad w_2=(-s,0,t),\qquad w_3=(0,-s,r).
\]
Any non-null candidate suffices. If all three are null, then
\[
 Q(w_1)+Q(w_2)-Q(w_3)=2t^2=0.
\]
Since \(2\ne0\), this gives \(t=0\). The vector \(w=(1,0,0)\) is then
perpendicular to \(v\) and has quadrance one. The construction includes
\(v=0\), and each branch uses just a scalar equality test.
\end{proof}

For a cap \(\rho=c(A_v)\) and such a \(w\), quadratic reduction gives
\[
 w\rho=\frac{(1-Q(v))w+2(v\times w)}{1+Q(v)}.
\]
The cross term is perpendicular to \(w\), and the rotation preserves its
quadrance. Therefore
\[
 s(w,w\rho)
 =1-\frac{(1-Q(v))^2}{(1+Q(v))^2}
 =4p(\rho)(1-p(\rho)).
\]
For a half turn the perpendicular vector is negated and the spread is zero.
This proves independence of the chosen admissible direction in all charts.
The proposed source repair is simply to replace the final existence sentence
of Theorem~33 by the four-candidate construction.

\section{Problem 80: the associative coefficient products}

The revised question concerns
\begin{equation}
 v\star u=\frac{v+u+\gamma(v\times u)}{1+\delta(v\cdot u)},
 \qquad\gamma\ne0,
 \label{eq:coefficient}
\end{equation}
on its displayed nonzero-denominator domain. ``Associative wherever
defined'' means equality of the two bracketings whenever all four products
occurring in them are defined. It does not mean that the two bracketings
have identical domains.

\begin{theorem}[Answer to Problem 80]
For \(\gamma,\delta\in F\) with \(\gamma\ne0\), the operation
\eqref{eq:coefficient} is associative wherever defined if and only if
\(\delta=-\gamma^2\).
\end{theorem}
\begin{proof}
Let \(e_1=(1,0,0)\), \(e_2=(0,1,0)\). For any nonzero \(a\)
with \(D=1+\delta a^2\ne0\), all four products for the triple
\((ae_1,ae_1,e_2)\) are then defined. The \(e_2\)-coordinate of
\((ae_1\star ae_1)\star e_2\) is one, whereas the same coordinate of
\(ae_1\star(ae_1\star e_2)\) is
\[
 \frac{1-\gamma^2a^2}{1+\delta a^2}.
\]
Equality gives \((\delta+\gamma^2)a^2=0\), hence the claimed relation.
The necessary probe or coefficient relation is obtained finitely. Try
\(a=1\). If its denominator vanishes, then \(\delta=-1\). If
\(\gamma^2=1\), the desired relation already follows. Otherwise take
\(a=\gamma\ne0\), whose denominator \(1-\gamma^2\) is nonzero. This
argument also covers characteristic three; no appeal to an infinite field
is needed.

For sufficiency assume \(\delta=-\gamma^2\). The invertible dilation
\(D_\gamma(v)=\gamma v\) transports \(\star\), including its denominator,
to the opposite anchor product:
\[
 D_\gamma(v\star u)=
 \frac{D_\gamma(v)+D_\gamma(u)
       +D_\gamma(v)\times D_\gamma(u)}
      {1-D_\gamma(v)\cdot D_\gamma(u)}.
\]
Its homogeneous product is the opposite of the associative quaternion
product \eqref{eq:quaternion}. Comparing affine representatives proves
associativity on every common domain, even without assuming that all
intermediate vectors are cap anchors.
\end{proof}

The fixed linear coefficients already make zero a two-sided identity. The
nonzero cross coefficient in the revised question permits the two-probe
argument above over every stated field. The retained formal theorem states
the stronger homogeneous rational classification: for
\[
 \frac{\alpha v+\beta u+\gamma(v\times u)}
      {\epsilon+\delta(v\cdot u)},
\]
zero is a two-sided identity and the product is associative wherever defined
exactly when
\[
 \epsilon\ne0,\qquad\alpha=\beta=\epsilon,
 \qquad\epsilon\delta=-\gamma^2.
\]
Thus the revised normalized problem has a positive classification answer;
the degenerate projection operations excluded by its new statement play no
role in that answer.

\section{Problem 81: finite fields}

Let \(F=\F_q\) with \(q\) odd, and write \(\eta(c)=1,-1,0\) according as
\(c\) is a nonzero square, a nonsquare, or zero. The constructions use a
supplied complete finite carrier and its exact square table.

\subsection{Vector counts and half-spread fibers}

\begin{lemma}[Ternary quadric count]
For \(a\in F\),
\[
 \#\{v:Q(v)=a\}=q^2+q\eta(-a).
\]
\end{lemma}
\begin{proof}
Put \(\varepsilon=\eta(-1)\), and let
\(C(b)=\#\{(r,s):r^2+s^2=b\}\). If \(-1=i^2\), the substitution
\((r,s)\mapsto(r+is,r-is)\) counts solutions of \(xy=b\), giving
\(C(0)=2q-1\) and \(C(b)=q-1\) for \(b\ne0\).
If \(-1\) is nonsquare then \(C(0)=1\). For \(b\ne0\), the two square
sets of size \((q+1)/2\) give a point on the conic
\(r^2+s^2=bt^2\). Lines through it parametrize \(q+1\) projective points;
none is at infinity, so \(C(b)=q+1\). In both cases
\[
 C(0)=q+(q-1)\varepsilon,\qquad C(b)=q-\varepsilon\quad(b\ne0).
\]
Sum \(C(a-t^2)\) over \(t\). There are one, two or zero exceptional
summands according as \(a\) is zero, a nonzero square or a nonsquare. These
three counts give the stated formula.
\end{proof}

\begin{theorem}[Rotation count and half-spread distribution]
The complete rotation count is \(q(q^2-1)\). Its half-spread fibers are
\[
 \#\{\rho:p(\rho)=h\}=
 \begin{cases}
 q^2,&h=0\text{ or }h=1,\\
 q^2+q\eta(-h(1-h)),&h\ne0,1.
 \end{cases}
\]
The two spinor-norm fibers each have \(q(q^2-1)/2\) elements.
\end{theorem}
\begin{proof}
In the affine chart \((1:v)\), exclude \(Q(v)=-1\), which removes
\(q^2+q\) vectors. At infinity the non-null lines number
\[
 \frac{(q^3-1)-(q^2-1)}{q-1}=q^2.
\]
Their sum is \(q^3-q\). For \(h\ne1\), the unique affine representative
satisfies \(Q(v)=h/(1-h)\). Apply the lemma, observing that the square
class of \(-h/(1-h)\) equals that of \(-h(1-h)\).
Half turns are precisely the \(q^2\) points at infinity.
Every nonzero field element is a sum of two squares, so both norm classes
have quaternion representatives. Multiplication by a representative of
nonsquare norm bijects their rotation fibers. They consequently have equal
size.
\end{proof}

For \(h\ne1\), every rotation in a half-spread fiber has spinor norm
\([1-h]\). For \(h=1\), its norm is the square class \([Q(v)]\) of its
axis. More explicitly, the half-turn axes of either prescribed square class
\([c]\) number
\[
 \frac{q^2+q\eta(-c)}2.
\]
Indeed, there are \((q-1)/2\) nonzero quadrances in that class, each having
the lemma's vector count, and each projective line contains \(q-1\)
nonzero vectors.

\subsection{Conjugacy types and parabolics}

The reciprocal factor of the characteristic polynomial is
\[
 X^2-2(1-2h)X+1,\qquad \Delta=-16h(1-h).
\]
For \(h\ne0,1\), it has distinct roots. Square \(-h(1-h)\) gives split
semisimple type; nonsquare \(-h(1-h)\) gives nonsplit semisimple type.
Through the retained split-quaternion matrix model, there is one projective
conjugacy class for each such fixed \(h\). Indeed, the matrix trace is nonzero
when \(h\ne1\); scale it to one. The identity
\(1-h=(\tr M)^2/(4\det M)\) then fixes its determinant at
\(1/(4(1-h))\), so both matrices have the same characteristic polynomial.
A non-scalar two-by-two matrix has a cyclic vector: test the two coordinate
vectors and their sum. If none were cyclic, the first two tests would make
the matrix diagonal and the third would force equal diagonal entries.
In the basis consisting of a cyclic vector and its image, the matrix is
the companion matrix of its characteristic polynomial. The two matrices
are therefore similar. A half turn instead has a non-scalar trace-zero
representative. Representatives with the same determinant square class can
be scaled to the same determinant and the same cyclic-basis argument
applies. The two determinant classes thus give its two conjugacy classes.

The zero fiber has characteristic polynomial \((X-1)^3\). Its identity
has anchor zero; the other \(q^2-1\) rotations have nonzero null anchors.
There are \(q+1\) null axis lines and \(q-1\) nonzero parameters over each.
These nonidentity parabolics form one conjugacy class in the \emph{full}
group \(\SO(3,F)\simeq\operatorname{PGL}_2(F)\). To see this, a nonidentity
unipotent projective matrix can be put in the form
\(\left(\begin{smallmatrix}1&a\\0&1\end{smallmatrix}\right)\),
\(a\ne0\); diagonal conjugation changes any such \(a\) to any other.
This statement is about the full group, not conjugacy restricted to its
index-two unit-representable subgroup.

\subsection{The distribution on one non-null axis}

Fix \(v\) with \(Q(v)\ne0\). The source's condition
\(Q(v)h(1-h)\) square for \(h\ne0,1\) holds for exactly
\begin{equation}
 \frac{q-2-\eta(-Q(v))}{2}
 \label{eq:axiscount}
\end{equation}
values. This is the count of points with \(y\ne0\) on
\[
 (h-\tfrac12)^2+Q(v)y^2=\tfrac14,
\]
after removing its two points with \(y=0\) and pairing the remaining
points as \(\pm y\). Thus the count is \((q-3)/2\) if \(-Q(v)\) is a
square, and \((q-1)/2\) otherwise. Each amount has exactly two inverse cap
rotations. The forbidden affine scales are counted by
\(\lambda^2Q(v)=-1\), equivalently the square-root fiber of \(-Q(v)\).
These are actual filtered-carrier and root-fiber counts in the formalisation,
including the literal examples over \(\F_3\) and \(\F_5\).

As a check on the global distribution, the half-spread counts for
\(\F_3,\F_5,\F_7\), in the order of the field's integer representatives,
are respectively
\[
 (9,9,6),\quad(25,25,20,30,20),\quad
 (49,49,56,42,42,42,56).
\]
They sum to \(24,120,336\). Over \(\F_5\) there are additionally
\(36\) null quaternion proportions, \(6\) null axis lines and \(24\)
nonidentity parabolics, agreeing with Example~70.

\section{Problem 82: what the spinor norm determines}

\subsection{Normalization and the kernel}

For a rational nonzero quaternion representative \(P\), \(N(P)>0\), and
\[
 \chi(\rho(P))=[N(P)].
\]
A rational norm-one representative exists exactly when this class is
trivial: \(N(P)=r^2\ne0\) permits division by \(r\), while any unit
rescaling exhibits such a square. The rotations with this property form
\(K=\ker\chi\), closed under products and inverses. For a half turn the
criterion is exactly that its supplied axis quadrance is a square.

\subsection{Finite generator lists and signed words}

\begin{theorem}[Complete finite spinor image]
Given a finite list of rational rotations \(g_1,\ldots,g_m\), one can
compute their complete finite spinor image, a word realizing every image
element, and exact kernel and kernel-coset tests for every supplied finite
signed word.
\end{theorem}
\begin{proof}
For positive \(c=a/b\) in reduced form, let \(S(c)\) be the finite set of
primes occurring to odd exponent in \(ab\). Finite factorization gives
\[
 [c]=[d]\ \Longleftrightarrow\ S(c)=S(d),
 \qquad S(cd)=S(c)\mathbin\triangle S(d).
\]
Choose retained representatives for the generators and take the union \(S\)
of their supports. Their parity vectors \(s_i\in\F_2^S\) span a finite
space \(V\). The spinor vector of a signed word is the sum of its generator
vectors; inverses have the same vector. Thus a word is in \(K\) exactly
when its vector is zero, and two words lie in the same \(K\)-coset exactly
when their vectors agree. Finite binary elimination computes the rank \(r\)
and the \(2^r\) image vectors, retaining generator words realizing them.
The quaternion product separately evaluates each word's exact rotation.
\end{proof}

This classifies the image and the quotient detected by \(\chi\), not the
generated subgroup internally. Different words can have the same spinor
vector and different rotation matrices. Matrix comparison decides any
supplied pair, but does not provide an isomorphism type, presentation or
complete relation theory of the generated subgroup.

\subsection{Finite-order values}

\begin{theorem}[Rational finite-order table]
If \(\rho\in\SO(3,\Q)\) is supplied with an exact positive order, then
the following table is exhaustive:
\[
\begin{array}{c|c|c}
\text{order}&p(\rho)&\chi(\rho)\\\hline
1&0&[1]\\
2&1&[Q(a)]\text{ for its nonzero axis }a\\
3&3/4&[1]\\
4&1/2&[2]\\
6&1/4&[3].
\end{array}
\]
\end{theorem}
\begin{proof}
Set \(\tau=\tr\rho-1\) and use the finite recurrence
\[
 L_0=2,\quad L_1=\tau,\quad L_{k+2}=\tau L_{k+1}-L_k.
\]
Cayley--Hamilton gives \(\tr(\rho^k)=1+L_k\). An identity power implies
\(L_n=2\), and the monic integer equation \(L_n(\tau)-2=0\) forces a
reduced rational \(\tau\) to be integral by denominator clearing and
B\'ezout. The Euclidean half-spread bounds give \(-2\le\tau\le2\).
The five integer values produce exactly the orders in the table; the
retained smaller nonidentity powers certify minimality. Finally
\(p=(2-\tau)/4\), and \(\chi=[1-p]\) away from half turns, give its last
two columns. At a half turn use \((0:a)\) instead.
\end{proof}

The order-two row is a complete answer for a supplied axis. A separate global
classification of the positive square classes represented by three rational
squares has not been supplied. No claim that every positive square class
occurs as an order-two value is made here.

\subsection{Retained finite-precision approximation}

The earlier development also proves a precise finite approximation statement:
for each rational rotation \(\rho\) and positive rational \(\varepsilon\),
it constructs \(\eta\in K\) with
\[
 \sum_{i,j}(\eta_{ij}-\rho_{ij})^2<\varepsilon.
\]
Choose the projectively equivalent sign of \(P=(x:v)\) with \(x\ge0\).
Put \(n=N(P)>0\), choose a sufficiently tight positive rational upper
square-root bracket endpoint \(u\), and set
\[
 A=u+x,\quad R=(A^2-Q(v):2Av),\quad D=A^2+Q(v),\quad U=D^{-1}R.
\]
The retained construction ensures \(D\ne0\). Exact expansion gives
\(N(R)=D^2\), hence \(N(U)=1\), and the rotation of \(U\) satisfies
\[
 \sum_{i,j}(\eta_{ij}-\rho_{ij})^2
 =\frac{8(u^2-n)^2Q(v)}{D^2n}.
\]
Finite integer square-root refinement makes this less than the requested
bound. This is a theorem for each finite request. It neither constructs a
completed topological space nor asserts a formal closure theorem.

\section{Problem 83: arbitrary nondegenerate ternary forms}

Let \(M\) be a supplied symmetric nonsingular three-by-three matrix over
\(F\), and put
\[
 B_M(u,v)=uMv^T,\quad Q_M(u)=B_M(u,u),\quad d=\det M,
 \quad C_M(u,v)=(u\times v)\adj(M).
\]
The cross on the right is the ordinary coordinate cross product. It is
essential to retain the determinant factors: a bare replacement of the
Euclidean dot product would give the wrong algebra for a general form.

\begin{theorem}[Form-relative projective Clifford correspondence]
The product and norm
\[
 (x:u)(y:v)=(xy-dB_M(u,v):xv+yu-C_M(u,v)),
 \qquad N_M(x:u)=x^2+dQ_M(u)
\]
are associative and multiplicative, with
\[
 P\bar P=\bar PP=(N_M(P):0),\qquad\overline{(x:u)}=(x:-u).
\]
For \(N_M(P)\ne0\), normalized conjugation
\[
 (0:w)\longmapsto N_M(P)^{-1}\bar P(0:w)P
\]
is a member of \(\SO(M,F)\). Every member occurs, and two non-null
representatives give the same rotation if and only if they differ by a
nonzero scalar.
\end{theorem}
\begin{proof}
Adjugate identities give the determinant-weighted Lagrange,
Binet--Cauchy and vector triple-product identities. Expanding the displayed
finite coordinates proves associativity, the conjugation identities and
norm multiplicativity. The same identities prove preservation of \(B_M\)
and determinant one of the action.

For surjectivity, the retained constructive two-reflection theorem gives
two non-null normals for any supplied rotation. Their pure Clifford
representatives give half turns; multiplying the two half turns gives the
same rotation as multiplying the reflections. This constructs a non-null
preimage, including for the identity. The construction selects non-null
vectors and displacements from explicit finite coordinate candidates.
For faithfulness, the action identities force a trivially acting non-null
representative to have vector part zero. Applying this to the relative
product of two representatives and cancelling the non-null factor gives the
nonzero scalar witness. Conversely a rescaling multiplies both action
numerator and norm by the same square.
\end{proof}

The affine chart \((1:u)\) has
\[
 p=\frac{dQ_M(u)}{1+dQ_M(u)},\qquad
 u\star_Mv=\frac{u+v-C_M(u,v)}{1-dB_M(u,v)}.
\]
For defined input caps and nonzero product denominator,
\[
 1+dQ_M(u\star_Mv)=
 \frac{(1+dQ_M(u))(1+dQ_M(v))}{(1-dB_M(u,v))^2}.
\]
Thus all the affine exceptional loci and their projective resolution extend
with the same retained form. No diagonalizing basis or square-root
normalization is required. The form \(I\) gives the Euclidean algebra;
\(\operatorname{diag}(1,-1,-1)\), with the source's explicit vector sign
conversion, gives the dihedron convention of the Minkowski companion
\cite{companion}.

\section{Rays, normalization and the size boundary}

A quaternion ray identifies \(P\) with \(aP\) only for \(a>0\).
A projective proportion permits every \(a\ne0\), so it consists of two
opposite rays. Quaternion multiplication and inversion respect positive
scaling; over an ordered field \(P^{-1}=\bar P/N(P)\) is in the same ray
as \(\bar P\), because \(N(P)>0\). Equality of rotation matrices means
the representatives are projectively equivalent, hence either on the same
ray or on opposite rays. Each rotation fiber therefore contains exactly two
distinct rays; the kernel consists of the scalar rays of \(1\) and \(-1\).

These quaternion rays must not be confused with rays of the \emph{anchor
vector}. The quaternions \(P\) and \(-P\) give the same rotation. The
nonzero affine anchors \(v\) and \(-v\) give inverse rotations, and are
distinct vector rays. The source's orientation data concerns this second
distinction.

A rational quaternion ray has a unit representative precisely when its norm
class is trivial. When it exists, its unit representative is unique within
that positive ray. Thus the full rational ray cover is strictly larger than
the rational norm-one cover. The Pauli matrix description and its algebraic
determinant-one unitary image are checked using a formal imaginary scalar;
no completed complex field or Lie group is needed for those identities.
The real spin-group and trigonometric identifications lie outside the present
formal constructive scope.

The infinitude clause of Theorem~74 and the infinite-index clause of
Corollary~75 have the following checked finite content. For every supplied
natural number \(n\), an executable construction returns \(n\) pairwise
inequivalent positive rational square classes and rotation preimages, and
\(n\) rotations pairwise separated modulo \(K\). Successive finite prime
constructions supply the classes; a four-square construction supplies exact
quaternion norms. The pairwise certificates are retained. This gives
arbitrary finite size without constructing an infinite quotient carrier.

An additional retained result, no longer a numbered open problem in the
revision, is point--plane polarity: \([x]\) corresponds to the plane with
normal \(x\). The identity
\((x\rho)\cdot(y\rho)=x\cdot y\) makes this correspondence equivariant
and preserves incidence. For non-null representatives, projective point
quadrance equals spread of the corresponding plane normals. Null
representatives retain incidence but do not admit that divided spread.

\section{Formal evidence and remaining questions}

\begingroup\raggedright
The reusable mathematical owners remain in \texttt{Purelib}; the publication
facade is
\path{formalisations/RationalRotations/RationalRotations.lean}. Its
claim ledger is \path{formalisations/RationalRotations/coverage.md}, with
the earlier draft's numbering preserved separately in
\path{formalisations/RationalRotations/coverage-2026-08-31.md}.
The following map gives direct entry points rather than copied proofs.
\par\endgroup

\begingroup\small
\begin{longtable}{>{\raggedright\arraybackslash}p{0.27\textwidth}>{\raggedright\arraybackslash}p{0.66\textwidth}}
\toprule
Result & Mathematical owner / retained entry point \\
\midrule
\endhead
Perpendicular witness &
\path{Purelib/Geometry/Quadratic/Space3/Perpendicular.lean};
\texttt{perpendicularNonNullWitness}. \\
Coefficient classification &
\path{Purelib/Geometry/Rotations/Cayley/CoefficientProducts.lean};
\newline\texttt{rational\_hardenedAssociativityClassification}. \\
Finite strata and conjugacy &
\path{Purelib/Geometry/Rotations/FiniteField/Classification.lean};
the original finite-field presentation and its classification certificate. \\
Fixed-axis amounts &
\path{Purelib/Geometry/Rotations/Axis/Finite.lean};
\texttt{admissibleAmounts\_length}, cached roots and two-scale fibers. \\
Generated spinor image &
\path{Purelib/Geometry/Rotations/Rational/Generators/Classification.lean};
\texttt{GeneratedSubgroups.classify}. \\
Finite orders and approximation &
\path{Purelib/Geometry/Rotations/Rational/Order.lean};
\path{Purelib/Geometry/Rotations/Rational/Kernel/Approximation.lean}. \\
General ternary correspondence &
\path{Purelib/Geometry/Rotations/Ternary/Theory.lean};
\texttt{TernaryForm.generalRotationTheory}, together with
\texttt{projectiveCliffordCorrespondence}. \\
Ray cover &
\path{Purelib/Geometry/Rotations/Projective/RayCover.lean};
the original \texttt{RayProportion} product, inverse and fiber laws. \\
\bottomrule
\end{longtable}
\endgroup

The mathematical tree at commit \texttt{a6219f044} passed all twelve
completion gates on 11 October 2026: warning-fatal full build and lint,
the full configured-script axiom audit, independent nanoda checking and the
tooling regressions. Builds used 64 Lean threads. The normal proof-closure
audit accepted 136,660 project declarations and 1,536 selected canonical
declarations with no allowed axioms; nanoda checked all 42 compatible
module sets. This certifies the selected mathematical proof closure, not a
separate strict audit of every internal meta-provider helper. These papers
change documentation only; they do not introduce new Lean declarations.
The finite-field counts, stratum certificates and parabolic conjugators
have retained Lean proofs. The cyclic-matrix argument above additionally
deduces the full semisimple and half-turn conjugacy classes in this paper;
it is not claimed as a newly checked standalone Lean classification theorem.

Problem~80 is answered over every field of characteristic different from
two, and the finite-field
and general-form results above answer Problems~81 and~83 in the displayed
constructive sense. Problem~82 remains open as an internal classification of
natural generated subgroups. A global list of order-two spinor classes and
any completed topological statement would need additional results. Historical
priority and numerical-performance claims have not been inferred from the
formalisation.

\begin{thebibliography}{9}
\raggedright
\bibitem{source} N.~J.~Wildberger,
\emph{Rational rotations of three dimensional Euclidean space}, supplied
revision \texttt{RationalRotationsQuaternions2.pdf}, 33 pages.
Source SHA-256:
\url{3e03d449b480ab4b27b520e8059495bab2e62cdf3852b1ff33e02a38c7ad8c53}.
\bibitem{companion} \emph{Issues and Resolutions to Open Problems in Rational
Rotations of Three Dimensional Minkowski Space}, companion in
\path{formalisations/MinkowskiRotations/IssuesAndOpenProblems.pdf},
12 October 2026.
\end{thebibliography}
\end{document}
